\documentclass[10pt,a4paper]{amsart}

\pdfinfoomitdate=1
\pdftrailerid{}
\pdfsuppressptexinfo=15

\usepackage[T1]{fontenc}
\usepackage{lmodern}
\usepackage[margin=25mm]{geometry}
\usepackage[protrusion=true,expansion=false]{microtype}
\usepackage{amsmath,amssymb,mathtools}
\usepackage{enumitem}
\usepackage{xcolor}
\usepackage[numbers,sort&compress]{natbib}
\usepackage[colorlinks=true,linkcolor=blue!55!black,citecolor=blue!55!black,urlcolor=blue!55!black]{hyperref}
\usepackage[nameinlink,noabbrev]{cleveref}

\setlist{nosep,leftmargin=*}
\raggedbottom
\allowdisplaybreaks
\emergencystretch=2em

\newtheorem{theorem}{Theorem}[section]
\newtheorem{proposition}[theorem]{Proposition}
\newtheorem{lemma}[theorem]{Lemma}
\newtheorem{corollary}[theorem]{Corollary}
\theoremstyle{definition}
\newtheorem{definition}[theorem]{Definition}
\theoremstyle{remark}
\newtheorem{remark}[theorem]{Remark}

\newcommand{\En}{\mathcal E_n}
\newcommand{\depth}{\operatorname{depth}}
\newcommand{\im}{\operatorname{im}}
\newcommand{\BA}{\boldsymbol\beta}

\title[Whole-array recomputation of border arrays]{Whole-Array Recomputation of Border Arrays:\\
Exact Two-Cycles, Sharp Transients, and Factorial Fibres}
\author{Anonymous}
\date{}
\subjclass[2020]{37P25, 68R15, 05A05}
\keywords{border array, prefix function, whole-array recomputation,
finite dynamical system, inversion sequence, fibre}

\begin{document}

\begin{abstract}
Let $\BA(w)$ be the ordinary border array of a finite word.  We treat the
entire integer array $\BA(w)$ as the next word, recompute all of its borders,
and repeat on the inversion-sequence carrier
$\En=\{(e_0,\ldots,e_{n-1}):0\le e_i\le i\}$.  This whole-array
recomputation differs from following failure links inside one fixed table.
We prove that its one-step image is exactly the valid border arrays.  At
length $n\ge2$, the recurrent set consists of $n-1$ explicit cycles of exact
period two; length one has one fixed point.  An indexed three-state mismatch
amplifier gives the sharp maximum transient
$0,0,1,2n-4$ for $n=1,2,3,n\ge4$.  We also prove, for every target, the
one-step bound $(n-1)!$ and show that equality occurs only at the all-zero
table and $(0,1,0,\ldots,0)$ when $n\ge2$.  Border computation, validation,
construction, realization, and census are credited background.  The owner
search is bounded; no novelty, priority, or external-release claim is made.
\end{abstract}

\maketitle

\section{Whole-array recomputation and the subtraction boundary}

For a word $w=w_0\cdots w_{n-1}$ over any alphabet, let
\begin{equation}
 \beta_i(w)=\max\{k:0\le k\le i,\ 
 w_0\cdots w_{k-1}=w_{i-k+1}\cdots w_i\}.             \label{eq:border}
\end{equation}
Thus $\beta_i(w)$ is the longest proper-border length of the prefix ending at
$i$, and $\beta_0(w)=0$.  Write
$\BA(w)=(\beta_0(w),\ldots,\beta_{n-1}(w))$.
The Morris--Pratt and Knuth--Morris--Pratt mechanisms compute these data in
linear time \cite{KnuthMorrisPratt1977}.  Validation and realization of
border or failure arrays, their generation, and their enumeration have
dedicated algorithms and structural theories
\cite{FranekEtAl2002,DuvalLecroqLefebvre2009,GawrychowskiJezJez2014}.
All of that one-step machinery is background here.

The carrier in this paper is
\begin{equation}
 \En=\{e=(e_0,\ldots,e_{n-1}):0\le e_i\le i\}.        \label{eq:carrier}
\end{equation}
Regarding $e$ as an integer word, define the self-map
\begin{equation}
                    \Pi_n(e)=\BA(e).                  \label{eq:update}
\end{equation}
It is closed on $\En$ because every proper border of an $(i+1)$-letter
prefix has length at most $i$.

The term \emph{whole-array recomputation} prevents a semantic collision.
For a fixed word $w$, the usual failure map sends a positive border length
$k$ to $\beta_{k-1}(w)$; iterating that scalar map follows the nested borders
of the same fixed word.  In \eqref{eq:update}, by contrast, the complete table
$\BA(e)$ becomes a new word:
\[
        e\longmapsto\BA(e)\longmapsto\BA(\BA(e))\longmapsto\cdots.
\]
No unqualified ``iteration of the prefix function'' is meant below.

A state is \emph{recurrent} if it belongs to a directed cycle of $\Pi_n$.
Its depth is
\[
 \depth(e)=\min\{t\ge0:\Pi_n^t(e)\text{ is recurrent}\}.
\]

\section{The one-step image and canonical two-cycles}

Call an integer array \emph{valid} if it is the border array of some word.

\begin{proposition}[exact image]\label{prop:image}
The image $\im\Pi_n$ is exactly the valid border arrays of length $n$.
\end{proposition}

\begin{proof}
Every output of $\Pi_n$ is valid by definition.  Conversely, let $p=\BA(w)$
be valid.  Replace the first distinct letter of $w$ by zero, the next
previously unseen letter by one, and so on.  The standardized letter at
position $i$ is at most $i$, so the standardized word $\widehat w$ belongs
to $\En$.  Standardization preserves every equality and inequality between
positions.  Borders depend only on this equality pattern; hence
$\BA(\widehat w)=\BA(w)=p$.
\end{proof}

For $1\le r<n$, define
\begin{align}
 A_r&=(0,1,\ldots,r,0,\ldots,0),\label{eq:A}\\
 B_{r+1}&=(\underbrace{0,\ldots,0}_{r+1},1,\ldots,1).\label{eq:B}
\end{align}

\begin{proposition}[canonical pairs]\label{prop:pairs}
For $1\le r<n$,
\begin{equation}
                  \Pi_n(A_r)=B_{r+1},\qquad
                  \Pi_n(B_{r+1})=A_r.                 \label{eq:pairs}
\end{equation}
In particular, these states form $n-1$ distinct cycles of exact period two.
\end{proposition}

\begin{proof}
Along the distinct initial slope of $A_r$, no nonempty border exists, so the
first $r+1$ entries of $\BA(A_r)$ are zero.  Every later prefix ends in zero
and has the one-letter border $0$.  If a proper suffix of length at least two
were also a prefix, its first letter would force it to start at a tail zero;
its second letter would then be zero, contradicting the second prefix letter
$1$.  Thus $\BA(A_r)=B_{r+1}$.

Inside the initial zero run of $B_{r+1}$, the longest proper border is one
shorter than the prefix, producing $0,1,\ldots,r$.  Consider a prefix that
ends in the one run.  A nonempty border must have length greater than $r+1$
so that its prefix copy also ends in one.  Its suffix copy must start at a
zero, but any proper such start lies inside the initial zero run and gives
strictly fewer leading zeros than the prefix copy.  Hence no nonempty border
exists after the zero run, $\BA(B_{r+1})=A_r$, and the templates are distinct.
\end{proof}

The remaining task is to prove that no other recurrent state exists and to
measure how quickly every state enters one of these pairs.

\section{The indexed mismatch amplifier}

Every border array $p$ satisfies the unit-growth inequality
\begin{equation}
                         p_i\le p_{i-1}+1.             \label{eq:growth}
\end{equation}
For a valid $p$ of length at least two, choose its canonical template as
follows.  If $p_1=1$, let $r$ be maximal such that
$p_0p_1\cdots p_r=01\cdots r$ and put $Q(p)=A_r$.  If the slope stops, a
realizing word has begun with $r+1$ equal letters.  Its next letter either
continues that equality and the slope, or differs and destroys every
positive border; maximality therefore forces $p_{r+1}=0$.

If $p_1=0$, let $k$ be the length of its maximal initial zero run and put
$Q(p)=B_k$.  If the run stops, \eqref{eq:growth} forces $p_k=1$.  Therefore
every noncanonical valid table agrees with its template through at least
three coordinates.  Define its first mismatch index by
\begin{equation}
 L(p)=\min\{i:p_i\ne Q(p)_i\},                         \label{eq:L}
\end{equation}
and put $L(p)=n$ when $p=Q(p)$.  The number $L(p)$ is also the length of the
agreement prefix.

\begin{lemma}[indexed mismatch transitions]\label{lem:mismatch}
Let $p$ be valid, $q=\Pi_n(p)$, and $L=L(p)<n$.
\begin{enumerate}[label=(\roman*)]
\item If $Q(p)=A_r$, then $p_L=1$, $Q(q)=B_{r+1}$,
      $L(q)=L$, and $q_L=2$.
\item If $Q(p)=B_k$, then $p_L\in\{0,2\}$ and $Q(q)=A_{k-1}$.
      For $p_L=0$, one has $L(q)=L$ and $q_L=1$; for $p_L=2$,
      one has $L(q)>L$.
\item After the extension in (ii), if another mismatch exists, it is an
      $A$-type mismatch whose actual value is one.
\end{enumerate}
Equivalently, the first-mismatch states obey
\begin{equation}
 A1\longrightarrow B2\longrightarrow\text{extension},\qquad
 B0\longrightarrow A1\longrightarrow B2\longrightarrow\text{extension}.
                                                               \label{eq:automaton}
\end{equation}
\end{lemma}

\begin{proof}
The border array of a prefix depends only on that prefix.  Thus $q$ agrees
with $\Pi_n(Q(p))$, the partner in \eqref{eq:pairs}, before index $L$.
If $Q(p)=A_r$, then $L$ lies after the first template-tail zero, so this
agreement contains the $r+1$ initial zeros of $B_{r+1}$ and its first
following one; hence $Q(q)=B_{r+1}$.  If $Q(p)=B_k$, then $L$ lies after the
first one following the zero run, so the shared prefix contains the complete
initial slope of $A_{k-1}$ and its first following zero; hence
$Q(q)=A_{k-1}$.  It is this complete shared prefix, not the second coordinate
alone, that fixes the partner parameter.

Suppose first that $Q(p)=A_r$.  The mismatch occurs after the first zero
following the slope, so $p_{L-1}=0$.  Its template value is zero; validity,
\eqref{eq:growth}, and mismatch force $p_L=1$.  As an integer word, the
prefix $p_0\cdots p_{L-1}$ has border length one.  The new letter $p_L=1$
matches $p_1=1$, so the border extends to length two.  Hence $q_L=2$ while
the $B_{r+1}$ template has value one, proving (i).

Now suppose $Q(p)=B_k$.  Here $p_{L-1}=1$ and the template value at $L$ is
one.  Inequality \eqref{eq:growth} leaves precisely $p_L=0$ or $p_L=2$.
The integer-word prefix $p_0\cdots p_{L-1}$ has border length zero.  If
$p_L=0$, it matches $p_0=0$ and creates border one, which is the $A1$ state.
If $p_L=2$, it does not match $p_0$, so $q_L=0$, the required $A_{k-1}$
value; agreement strictly extends.  Any later mismatch is now in an
$A$-template tail after a zero.  The preceding $A_r$-case argument forces its
actual value to be one.  This proves (ii), (iii), and \eqref{eq:automaton}.
\end{proof}

\begin{theorem}[recurrent atlas and upper clock]\label{thm:upper}
At length one, $(0)$ is the unique recurrent state and is fixed.  For
$n\ge2$, the recurrent set is exactly
\[
                    \{A_r,B_{r+1}:1\le r<n\},
\]
so it consists of the $n-1$ exact two-cycles in \cref{prop:pairs}.  For
$n\ge4$,
\begin{equation}
 \max_{p\in\im\Pi_n}\depth(p)\le2n-5,\qquad
 \max_{e\in\En}\depth(e)\le2n-4.                    \label{eq:upper}
\end{equation}
\end{theorem}

\begin{proof}
A noncanonical valid table has $L\ge3$.  By
\cref{lem:mismatch}, its first mismatch is passed in at most three updates.
After that extension, each later mismatch is passed in at most two updates.
Since every extension increases $L$ by at least one, the number of updates
before $L=n$ is at most
\[
              3+2(n-L-1)\le3+2(n-4)=2n-5.            \label{eq:count}
\]
At $L=n$ the table is one of the canonical templates and is recurrent.
Strict growth of $L$ excludes recurrence before that time.  Every state of
$\En$ enters the valid image after one update by \cref{prop:image}, giving
the second bound.  The length-one statement and \cref{prop:pairs} finish the
recurrent classification.
\end{proof}

\section{Sharp trajectories and the small boundaries}

For $n\ge4$, put
\begin{align}
 p_n&=(0,0,1,0^{n-3}),& e_n&=(0,1,0,2,1^{n-4}),       \label{eq:witness}\\
 X_j&=(0,0,1^j,2,0^{n-j-3}) &&(1\le j\le n-3),       \label{eq:X}\\
 Y_j&=(0,1,0^{j+1},1,2^{n-j-4}) &&(0\le j\le n-4). \label{eq:Y}
\end{align}

\begin{proposition}[sharp witness trajectory]\label{prop:witness}
The following identities hold:
\begin{align}
 \Pi_n(e_n)&=p_n,& \Pi_n(p_n)&=Y_0,\label{eq:first}\\
 \Pi_n(Y_j)&=X_{j+1} &&(0\le j\le n-4),\label{eq:YtoX}\\
 \Pi_n(X_j)&=Y_j &&(1\le j\le n-4),\label{eq:XtoY}\\
 \Pi_n(X_{n-3})&=A_1.                                \label{eq:endpoint}
\end{align}
Consequently $\depth(p_n)=2n-5$ and $\depth(e_n)=2n-4$.
\end{proposition}

\begin{proof}
All identities follow from equality blocks, which we spell out to fix the
endpoints.  In $X_j$, the initial $00$ creates border values $0,1$; the
following $j$ ones and the letter $2$ have border zero.  If a zero tail is
present, its first letter has border one and every later zero has border two.
This is $Y_j$.  When $j=n-3$, there is no zero tail, and the output is
$A_1=(0,1,0^{n-2})$.

In $Y_j$, the second letter differs from the initial zero, giving border
zero.  Each of the next $j+1$ zeros has the one-letter border $0$ and no
longer border.  The following one completes the border $01$ of length two.
For a prefix ending in the trailing $2$-run, a proper border would contain
fewer terminal twos in the initial copy than in the suffix; equality of those
counts occurs only for the full, nonproper prefix.  Thus every trailing $2$
has border zero.  The result is $X_{j+1}$.  The same inspection gives the two
initial identities in \eqref{eq:first}.

Thus the orbit of $p_n$ visits
$Y_0,X_1,Y_1,X_2,\ldots,X_{n-3},A_1$ at successive times and first reaches
the canonical two-cycle after $2n-5$ updates.  None of the preceding states
is canonical.  Since $e_n$ maps to $p_n$, it has one additional transient
step.
\end{proof}

\begin{theorem}[piecewise sharp maximum]\label{thm:sharp}
For every $n\ge1$,
\begin{equation}
 \max_{e\in\En}\depth(e)=
 \begin{cases}
 0,&n=1,\\
 0,&n=2,\\
 1,&n=3,\\
 2n-4,&n\ge4.
 \end{cases}                                           \label{eq:sharp}
\end{equation}
\end{theorem}

\begin{proof}
For $n\ge4$, \cref{thm:upper,prop:witness} give matching bounds.  At length
one there is only the fixed state.  At length two the two carrier states are
$A_1$ and $B_2$, hence both are recurrent.  At length three, the four
templates are recurrent, while the remaining two states $(0,0,2)$ and
$(0,1,1)$ enter a template in one update.  This proves every boundary in
\eqref{eq:sharp}.
\end{proof}

\section{Every target fibre and its two unique maxima}

The next theorem applies to every target $p\in\En$, whether or not $p$ is a
valid border array.

\begin{theorem}[factorial fibre extremum]\label{thm:fibre}
For every $p\in\En$,
\begin{equation}
                         |\Pi_n^{-1}(p)|\le(n-1)!.     \label{eq:fibre}
\end{equation}
When $n\ge2$, equality holds exactly for
\begin{equation}
                     p=0^n\quad\text{and}\quad
                     p=A_1=(0,1,0^{n-2}).             \label{eq:max-targets}
\end{equation}
At $n=1$, the sole fibre has size one.
\end{theorem}

\begin{proof}
Expose a source $e\in\En$ from left to right.  Its first letter is
$e_0=0$.  The prescribed value $p_1$ uniquely determines $e_1$: value one
requires $e_1=0$, and value zero requires $e_1=1$.  At position $i\ge2$, a
positive prescribed border $p_i=k$ forces the last source letter to match
the last letter of the corresponding prefix, namely $e_i=e_{k-1}$; it leaves
at most one choice.  If $p_i=0$, then at least $e_i\ne e_0=0$, leaving at
most the $i$ choices $1,\ldots,i$.  Multiplication proves \eqref{eq:fibre};
an invalid target simply has no sources.

Equality in the product requires $p_i=0$ at every $i\ge2$, because replacing
the factor $i\ge2$ by at most one is strict.  Since $p_1\in\{0,1\}$, only the
two targets in \eqref{eq:max-targets} remain.  We now prove that both attain
the bound without using the false general shortcut that a nonzero final
letter alone forbids every longer border.

When $n=2$, the product over $i=2,\ldots,n-1$ is empty and equals one: the
sources $(0,1)$ and $(0,0)$ map respectively to $0^2$ and $A_1$.  Hence
assume $n\ge3$ in the suffix arguments that follow.

For target $0^n$, choose $e_1=1$ and choose each later $e_i$ arbitrarily in
$\{1,\ldots,i\}$.  Position zero is the only zero.  Every proper suffix of a
prefix therefore starts with a nonzero letter and cannot equal a prefix,
which starts with zero.  All border values are zero, and there are
$\prod_{i=2}^{n-1}i=(n-1)!$ sources.

For target $A_1$, choose $e_1=0$ and again choose
$e_i\in\{1,\ldots,i\}$ for $i\ge2$.  The prefix of length two has border one.
For any longer prefix, a proper suffix starting at position at least two
fails in its first letter.  The only other possible proper suffix starts at
position one; its second letter is $e_2\ne0=e_1$, so it also fails.  Hence all
later border values are zero, giving another $(n-1)!$ sources.  These are the
only equality cases.  For $n=1$, both the carrier and its image consist of
$(0)$.
\end{proof}

\section{Exact control, scope, and limitations}

The paper-local verifier is deterministic, dependency-free, and uses exact
Python integers.  It exhausts all
\[
                  \sum_{n=1}^{9}n!=409{,}113
\]
carrier states and the same number of target cells.  Through length eight it
compares the linear border routine with literal longest-prefix and
longest-suffix comparison.  It checks every recurrent state and period, every
instance of \cref{lem:mismatch}, both depth bounds, the displayed sharp
trajectory, every target fibre, and both unique maximizers.  An independent
standardization check covers $3{,}279$ words, and the sharp witnesses are
continued through every length $4\le n\le32$.  The frozen run makes
$1{,}694{,}506$ exact assertions.  These computations search for
counterexamples; they do not prove the all-length statements.

The subexceedant carrier, ordinary longest-border convention, and synchronous
whole-array recomputation are essential.  Strict border arrays, optimized
failure functions, a bounded alphabet, or scalar failure-link descent define
other systems.  The carrier is an inversion-sequence encoding, so generic
permutation or factorial language receives no credit; likewise, a sharp
clock plus a fibre maximum is only a presentation silhouette.  The internal
collision controls P105, P112-C6, P122, and the owner-killed same-batch PR1
are therefore excluded from the claimed residual.

The literature search can establish a source hit but cannot certify an
absence.  The bounded search performed for this project did not locate the
whole-array recomputation system or the theorem package above; this non-hit
is not novelty or priority evidence.  Authorship, posting, submission,
specialist contact, and all external-release actions remain on hold.

\bibliographystyle{plainnat}
\bibliography{references}

\end{document}
